\subsection{极值元分解}

设 A 为整环，K 为其分式域，U 为 A 的可逆元素乘法群。回忆（《代数》第VI章，第1节，第5号），存在从 \(\mathbf{K}^* / \mathbf{U}\) 到 A 的非零分式主理想群 \(\mathcal{P}^*\) 的典范同构。定理1的条件 (b) 可翻译如下：

命题2. 设 A 为整环。A 为因子整环的充要条件是，存在 A 的子集 P，使 A 的每个 \(a \in A - \{0\}\) 都能唯一写成 \(a = u \prod_{p \in P} p^{n(p)}\)，其中 \(u \in U\)，且 \(n(p)\) 是除有限多个外为零的正整数。

若 P 满足此条件，显然它的所有元素都是极值元，且 A 的每个极值元都与 P 中唯一的一个元素相伴。回忆，P 此时称为 A 的极值元代表系（《代数》第VII章，第1节，第3号，定义2）。

始终假设 A 是因子整环。已经看到（第2节，定理1），群 \(\mathcal{P}^*\) 是一个格。因此可以应用《代数》第VI章第1节第9号和第13号的结果。特别地，\(\mathbf{K}^*\) 中的每个元素都可以以不可约分式的形式，按照本质唯一的方式写出。任意两个
元素 \(a, b\) 属于 \(\mathbf{K}^*\)，都有最大公因元和最小公倍元；若 \(a = u\prod_{p\in \mathbb{P}}p^{n(p)}\)，且

\[
b = u ^ {\prime} \prod_ {p \in P} p ^ {m (p)}
\]

这是 \(a\) 和 \(b\) 作为极值元乘积的分解。于是：

\begin{enumerate}
\def\labelenumi{(\arabic{enumi})}
\tightlist
\item
\end{enumerate}

\[
\mathbf {g . c . d .} (a, b) = w \prod_ {p \in P} p ^ {\inf (m (p), n (p))}\tag{2}
\]

\[
\mathbf {l . c . m .} (a, b) = w ^ {\prime} \prod_ {p \in P} p ^ {\sup (m (p), n (p))}
\]

其中 \(w, w'\) 属于 U。我们特别得到《代数》第VII章第1节第3号的结果。

对于所有 \(p \in P\)，映射 \(a \mapsto n(p)\) 是 K 上的离散赋值 \(v_{p}\)，其赋值环为 \(A_{Ap}\)。由定理1(e) 可知，\(v_{p}\) 正是 A 的本性赋值，而理想 \(A_{p}\)（\(p \in P\)）正是 A 的高度为1的素理想。

